A group that wants to detect a new finding on CT, with a few hundred to a few thousand labeled scans, can now start from one of several publicly released 3-D CT encoders, among them TAP-CT \citep{veenboer2025tapct}, SPECTRE \citep{claessens2025spectre}, Pillar-0 \citep{agrawal2025pillar0}, CT-FM \citep{pai2025ctfm}, Merlin \citep{blankemeier2026merlin}, CT-CLIP \citep{hamamci2026ctrate}, FM-HCT \citep{zhu2025fmct} and VoCo \citep{wu2024voco,wu2024vocolarge}; 2-D radiology models such as Curia process a volume slice by slice \citep{dancette2025curia}. Their evaluations rank representations: TAP-CT reports a linear layer on frozen tokens for segmentation and an attention-pooled head on frozen class tokens for classification \citep{veenboer2025tapct}, SPECTRE a nearest-neighbour classifier on frozen embeddings \citep{claessens2025spectre}, and CT-FM end-to-end fine-tuning \citep{pai2025ctfm}, while for FM-HCT probing and fine-tuning are compared on the same tasks \citep{zhu2025fmct}. However, such a table does not settle what the group has to decide before training: whether to keep the encoder frozen, adapt it with a low-rank update \citep{hu2022lora} or fine-tune part or all of it; whether to train one model per finding or one for all findings of a cohort; and, for segmentation, which head turns tokens into a voxel map. These choices change accuracy, compute and the number of models to maintain, and a frozen probe does not predict them: in our comparison, SPECTRE trails TAP-CT on the RSNA 2022 cervical-spine labels with the linear probe ($\numBbdiffRsnaBZBBLinMeandiff$ AUROC) but not once both are adapted with LoRA ($\numBbdiffRsnaBZBBLoraMeandiff$), and full fine-tuning with a standard setting collapsed on one of the encoders.

The evidence for these choices comes mostly from 2-D backbones and single tasks. On 2-D medical classification, LoRA matches or surpasses full fine-tuning of ImageNet-pretrained transformers \citep{dutt2024peft,zhu2024melo}, and probing and fine-tuning of DINOv2 \citep{oquab2023dinov2} have been compared on radiology tasks \citep{baharoon2023dinov2med}. For 3-D classification, \citet{liu2025anymc3d} found that frozen probing understates 2-D and native 3-D foundation models, and on one of our labels, RSNA 2023 spleen injury, full fine-tuning of a convolutional and a Swin 3-D model surpassed LoRA, which surpassed linear probing. Parameter-efficient adaptation of 3-D medical encoders has otherwise been studied mostly for segmentation: adapters that lift the 2-D Segment Anything model \citep{kirillov2023sam} to volumes \citep{gong20243dsamadapter,chen2024masam,wu2025medsa}, prompt tuning of a self-supervised 3-D U-Net \citep{fischer2024prompttuning}, LoRA, adapters and full fine-tuning of a supervised 3-D CT model in few-shot transfer \citep{silvarodriguez2025fseft}, and automated adapter configuration for few-shot 3-D segmentation \citep{ly2026seapeft}; for a 3-D CT vision--language model, full fine-tuning and parameter-efficient fine-tuning were close on one classification task \citep{kartika2026hcc}. Each of these studies uses one backbone and one task, and variants of LoRA that reparameterize or rescale the update \citep{liu2024dora,kalajdzievski2023rslora}, probing before fine-tuning \citep{kumar2022lpft} and fine-tuning only the attention weights \citep{touvron2022three} were not part of them. Joint training on many labels can help \citep{kim2023ctan} or hurt through conflicting gradients \citep{yu2020pcgrad,zhao2023mtlreview}, and task-specific LoRA weights generated by a hypernetwork have been proposed for chest CT \citep{liu2026hyperct}; the closest CT foundation-model studies do not compare joint and per-task adaptation of a 3-D encoder on the same labels \citep{veenboer2025tapct,liu2025anymc3d}. For segmentation, nnU-Net \citep{isensee2021nnunet} and models pretrained on segmentation labels \citep{li2024suprem,he2024vista3d} remain the references.

In this work, we compare what to train when adapting a pretrained 3-D CT encoder to new labels, under one protocol, on four public pretrained backbones of three architectural families: two plain vision transformers (TAP-CT, SPECTRE), a convolutional-stem multi-scale transformer (Pillar-0) and a hierarchical Swin transformer (VoCo). For classification we compare \numClsMethodsN{} methods, from frozen linear and attentive probes through low-rank adapters and their variants to attention-only and full fine-tuning, on \numLabelsN{} binary labels from five public cohorts. For segmentation we compare \numSegMethodsN{} methods, from a linear probe and a multi-depth decoder on the frozen encoder to adapters and full fine-tuning on top of the decoder, on \numStructuresN{} structures from five public datasets. Every method is trained once per label and once jointly for all labels of a cohort, with one default setting and budget, and compared at label level with intervals over cohorts, labels and held-out cases. Our contributions can be summarized as follows:
\begin{itemize}\setlength{\itemsep}{1pt}
  \item We provide a controlled comparison of \numClsMethodsN{} classification and \numSegMethodsN{} segmentation adaptation methods on four 3-D CT backbones (\numRunsMainN{} training runs, plus \numRunsRevisionN{} repeats and controls), trained per task and jointly, with per-case predictions released.
  \item We show which part of the encoder is worth training for each task and how far the answer transfers between backbone families, and we quantify the cost of one joint model against one model per label.
  \item We condense the results into a practical guide that separates consistent findings from conditional and uncertain ones.
\end{itemize}
